% ---------------------------------------------------------------------------
%  QML Lab research project 05: QAOA for MaxCut with Learned Initial Angles
%  Build:  latexmk -pdf main.tex   (needs qmlproject.sty, harvardqml.bst, references.bib)
% ---------------------------------------------------------------------------
\documentclass[11pt,a4paper]{article}
\usepackage{qmlproject}
\qmlheader{Project 05: QAOA learned angles}

\hypersetup{pdftitle={QML Lab project 05: QAOA for MaxCut with Learned Initial Angles}}

\begin{document}

\begin{center}
  {\small\textcolor{black!55}{QML Lab \textbullet\ research project 05}}\\[0.5em]
  {\LARGE\bfseries QAOA for MaxCut with Learned Initial Angles}\\[0.6em]
  {\large Predicting good QAOA parameters from graph features to cut optimisation cost.}\\[0.6em]
  {\small Living document \textbullet\ last updated \today}
\end{center}

\begin{projectcard}
\small
\begin{tabularx}{\linewidth}{@{}>{\bfseries}l X >{\bfseries}l X@{}}
Quadrant  & \QC\ (quantum data or quantum algorithm, classical learner) & Track & Variational \\
Difficulty & Medium & Tools & PennyLane, NetworkX, scikit-learn \\
Builds on & Modules 3 and 5 of the QML Lab & Workspace & \path{qml_projects/qaoa-maxcut-learned-init/} \\
\end{tabularx}
\end{projectcard}

\paragraph{Summary.}
The Quantum Approximate Optimization Algorithm (QAOA) solves combinatorial problems such as MaxCut
with a short parameterised circuit, but finding good angles costs many circuit evaluations. Optimal
angles are known to concentrate across similar graphs. This project trains a classical model that
maps simple graph features to good starting angles, and measures how many optimisation steps that
saves. The learner is classical and learns about a quantum algorithm, so the project sits on the
\QC\ side of the taxonomy.

% ===========================================================================
\section{Background and key concepts}
% ===========================================================================
\subsection{MaxCut}
For a graph $G=(V,E)$, MaxCut asks for a split of the vertices into two sets that cuts as many edges
as possible. With spins $z_i\in\{-1,+1\}$ the cost is
\begin{equation}
  C(z)=\sum_{(i,j)\in E}\frac{1-z_iz_j}{2}.
\end{equation}
The problem is NP-hard; the classical Goemans--Williamson algorithm guarantees a ratio of about
0.878 of the optimum \citep{goemans1995}.

\subsection{QAOA}
QAOA \citep{farhi2014} prepares
\begin{equation}
  \ket{\gamma,\beta}=\prod_{l=1}^{p}e^{-i\beta_lB}\,e^{-i\gamma_lC}\ket{+}^{\otimes n},
  \qquad B=\sum_jX_j,
\end{equation}
and maximises $F_p(\gamma,\beta)=\braket{\gamma,\beta|C|\gamma,\beta}$ with a classical optimiser. Quality
is measured by the approximation ratio $r=F_p/C_{\max}$. The landscape has many local optima, so
initial angles matter.

\subsection{Angle concentration and transfer}
\citet{brandao2018} showed that for fixed angles the QAOA objective concentrates over typical
instances of a graph family, so good angles for one graph are good for similar graphs.
\citet{zhou2020} found smooth patterns in optimal angles across depths and used them to build
initialisations for larger $p$. \citet{egger2021} warm-start QAOA from a classical relaxation. Learning
a map from graph features to angles generalises these ideas with data.

% ===========================================================================
\section{Research questions}
% ===========================================================================
The project is open-ended: these questions are the starting point, not the finish line. Add new
ones to the README and the journal as they come up.

\begin{rqlist}
  \item Which graph features (size, density, degree statistics, clustering) predict good angles?
  \item Do angles learned on 6--8-node graphs transfer to 10--12-node graphs?
  \item How much do predicted angles reduce optimiser steps compared with random and with fixed mean angles?
\end{rqlist}

% ===========================================================================
\section{Method}
% ===========================================================================
\subsection{Data}
Random graphs generated with NetworkX \citep{hagberg2008}: 3-regular graphs and Erd\H{o}s--R\'enyi
graphs with $p_{\text{edge}}\in\{0.3,0.5\}$, 6 to 12 nodes. For each graph, the best angles for depth
$p=1,2$ are found by optimisation from 20 random restarts; these are the regression targets. Angles
are mapped to a canonical range to remove symmetries before training.

\subsection{Models}
Regressors from graph features to $(\gamma,\beta)$: ridge regression, random forest and a small MLP.
QAOA itself is simulated exactly with PennyLane; optimisation of the angles uses Adam or COBYLA with a
fixed step budget.

\subsection{Baselines}
Random initial angles; fixed angles equal to the mean of the training targets (pure concentration);
the INTERP heuristic of \citet{zhou2020} for $p>1$; Goemans--Williamson for the cut quality reference.

\subsection{Metrics}
\begin{itemize}
  \item Approximation ratio at initialisation and after $K$ optimiser steps.
  \item Steps needed to reach 95\,\% of the best-known ratio.
  \item Angle prediction error on held-out graphs.
\end{itemize}

% ===========================================================================
\section{Planned experiments}
% ===========================================================================
Each experiment becomes one \texttt{[[experiments]]} entry in the workspace, so its runs, parameters
and metrics are stored and compared in the portal.

\begin{table}[h]
\centering\small
\begin{tabularx}{\linewidth}{@{}l>{\raggedright}X>{\raggedright}X>{\raggedright\arraybackslash}X@{}}
\toprule
ID & Experiment & What varies & Main metrics \\
\midrule
  E1 & Data set & graph family, $n=6..10$, $p=1$ & target quality \\
  E2 & Regressors & ridge, forest, MLP & init ratio, steps saved \\
  E3 & Size transfer & train $n\le8$, test $n=10,12$ & init ratio \\
  E4 & Deeper QAOA & $p=2,3$ & steps saved vs INTERP \\
  E5 & Graph families & regular vs Erd\H{o}s--R\'enyi & transfer across families \\
\bottomrule
\end{tabularx}
\caption{Experiment plan. Results go into the journal and the results log (\cref{sec:log}).}
\end{table}

% ===========================================================================
\section{Running it in the QML Lab}
% ===========================================================================
The workspace lives in \path{qml_projects/qaoa-maxcut-learned-init/}. Scripts import three helpers from
\path{qml_lab/plots.py}:
\begin{itemize}
  \item \texttt{params()} reads the parameters from the portal's run form (defaults when run by hand);
  \item \texttt{save()} stores figures with the run;
  \item \texttt{log\_metrics()} records numbers that appear in the run table.
\end{itemize}

\begin{lstlisting}
qml_projects/qaoa-maxcut-learned-init/
  project.toml          metadata + experiments
  README.md             goal, open questions, ideas
  experiments/          runnable scripts
  docs/main.tex         this document
\end{lstlisting}

First experiment: This experiment is \emph{planned}: write the script, add the block below to \texttt{project.toml}, then run \texttt{python manage.py sync\_projects}.

\begin{lstlisting}[language={}]
[[experiments]]
slug = "learned-init"
title = "Learned initial angles for depth-1 QAOA"
script = "experiments/e01_learned_init.py"
figure = "steps_saved.png"

[experiments.params]
n_nodes = 8
graph_type = "regular3"
depth_p = 1
n_graphs = 200
restarts = 20
seed = 0
\end{lstlisting}

\begin{lstlisting}[language=bash]
# from the portal folder: pick up changes to project.toml
python manage.py sync_projects
cd qml_projects/qaoa-maxcut-learned-init
# run by hand with the default parameters
PYTHONPATH=../../qml_lab python experiments/e01_learned_init.py
\end{lstlisting}

% ===========================================================================
\section{Quality bar}
% ===========================================================================
\begin{itemize}
  \item A classical baseline with a similar number of trainable parameters, tuned with the same care \citep{bowles2024}.
  \item At least five random seeds per configuration, reported as mean $\pm$ standard deviation.
  \item How results change with the number of qubits.
  \item At least one run under a realistic noise model.
  \item Report the classical Goemans--Williamson ratio next to every QAOA result.
\end{itemize}

% ===========================================================================
\section{Open directions}
% ===========================================================================
\begin{itemize}
  \item Graph neural networks as the angle predictor.
  \item Weighted MaxCut and other QUBO problems.
  \item Warm-starting from classical relaxations \citep{egger2021}.
  \item Run predicted-angle QAOA on hardware (Project 16).
\end{itemize}

% ===========================================================================
\section{Risks and limitations}
% ===========================================================================
\begin{itemize}
  \item Angle symmetries make targets ambiguous unless they are canonicalised.
  \item Small graphs are easy classically; the study is about optimisation cost, not about beating classical solvers.
  \item Generating targets needs many simulations; keep graphs small at first.
\end{itemize}

% ===========================================================================
\section{Results log}
\label{sec:log}
% ===========================================================================
Summarise findings here as they become stable; day-to-day notes belong in the portal journal.

\begin{table}[h]
\centering\small
\begin{tabularx}{\linewidth}{@{}l l >{\raggedright}X >{\raggedright\arraybackslash}X@{}}
\toprule
Date & Run & Finding & Next step \\
\midrule
 & & & \\
 & & & \\
 & & & \\
\bottomrule
\end{tabularx}
\end{table}

\bibliographystyle{harvardqml}
\bibliography{references}

\end{document}
